\documentclass[12pt,letterpaper]{article}



%%
%% Required packages:
%%
\usepackage{etex}
\usepackage{amsmath}
\usepackage{amssymb}
\usepackage{amstext}
\usepackage{amsthm}
\usepackage{fancyhdr,fancybox}
\usepackage{mathrsfs} % need for \mathscr command
\usepackage{multicol}
\usepackage[normalem]{ulem}
%\usepackage{pictex,pictexplus}
% \usepackage{pictexwd}
\usepackage{rotating}
\usepackage{url}
\usepackage{xfrac}
\usepackage{booktabs}
\usepackage{color,soul}
\usepackage{comment}

\usepackage{hyperref}
\hypersetup{
    colorlinks=true,     % false: boxed links; true: colored links
    linkcolor=blue,      % color of internal links
    citecolor=blue,      % color of links to bibliography
    filecolor=blue,      % color of file links
    urlcolor=blue        % color of external links
}


\usepackage{tikz}
\usetikzlibrary{backgrounds,calc}

%%
%% Common formatting commands
%%
\setlength{\parindent}{0in}
\setlength{\textwidth}{7in}
\setlength{\evensidemargin}{-0.25in}
\setlength{\oddsidemargin}{-0.25in}
\setlength{\parskip}{.5\baselineskip}
\setlength{\topmargin}{-0.5in}
\setlength{\textheight}{9in}

%               Problem and Part
\newcounter{problemnumber}
\newcounter{partnumber}
\newcounter{subpartnumber}
\newcommand{\Problem}{%
  \refstepcounter{problemnumber}%
  \setcounter{partnumber}{0}%
  \setcounter{subpartnumber}{0}%
  \item[\makebox{\hfil\textbf{\theproblemnumber.}\hfil}]%
  }
\newcommand{\InProblem}[2][1.6]{%
  \refstepcounter{problemnumber}%
  \setcounter{partnumber}{0}%
  \setcounter{subpartnumber}{0}%
  \textbf{\theproblemnumber.}\ \ \parbox[t]{#1in}{#2}%
}
\newcommand{\InSmallProblem}[1]{\InProblem[1.05]{#1}}
\newcommand{\NoProblem}[1]{%
  \mbox{\hfil\textbf{#1}\hfil}%
  }
\newcommand{\UnProblem}{%
  \refstepcounter{problemnumber}%
  \setcounter{partnumber}{0}%
  \setcounter{subpartnumber}{0}%
  \mbox{\hfil\textbf{\theproblemnumber.}\hfil}%
  }
\newcommand{\InFlexProblem}[1]{%
  \refstepcounter{problemnumber}%
  \setcounter{partnumber}{0}%
  \setcounter{subpartnumber}{0}%
  \textbf{\theproblemnumber.}\ \ #1%
}
%%  Part:
\newcommand{\Part}{%
  \renewcommand{\thepartnumber}{(\alph{partnumber})}%
  \refstepcounter{partnumber}
  \setcounter{subpartnumber}{0}%
  \item[(\alph{partnumber})]%
}
\newcommand{\InPart}[2][1.75]{%
  \renewcommand{\thepartnumber}{(\alph{partnumber})}%
  \refstepcounter{partnumber}%
  \setcounter{subpartnumber}{0}%
  (\alph{partnumber})\ \ \parbox[t]{#1in}{#2}%
}
\newcommand{\UnPart}{%
  \refstepcounter{partnumber}%
  \renewcommand{\thepartnumber}{(\alph{partnumber})}%
  \setcounter{subpartnumber}{0}%
  (\alph{partnumber})%
}
\newcommand{\InLargePart}[1]{\InPart[3.05]{#1}}
\newcommand{\InFullPart}[1]{\InPart[6.2]{#1}}
\newcommand{\InSmallPart}[1]{\InPart[1.05]{#1}}
%% SubPart:
\newcommand{\SubPart}{%
  \refstepcounter{subpartnumber}%
  \item[(\roman{subpartnumber})]%
  }
\newcommand{\InSubPart}[2][2.25]{%
  \stepcounter{subpartnumber}%
  (\roman{subpartnumber})\ \ \parbox[t]{#1in}{#2}%
}
\newcommand{\InSmallSubPart}[1]{\InSubPart[1.5]{#1}}
\newcommand{\InTinySubPart}[1]{%
  \stepcounter{subpartnumber}%
  (\roman{subpartnumber})\ \ \makebox{#1}%
}
\newcommand{\InMySubPart}[2][2.25]{%
  \stepcounter{subpartnumber}%
  \parbox[t]{#1in}{\makebox[0.3in][l]{(\roman{subpartnumber})}\ \ {#2}}%
}
\newcommand{\TrueFalse}{\textbf{T}\ \ \textbf{F}\ \ }
\newcommand{\TFPart}[1]{% 
  \stepcounter{partnumber}%
  \setcounter{subpartnumber}{0}%
  \item[(\alph{partnumber})] \TrueFalse\parbox[t]{5.5in}{#1}}

\newcommand{\Points}[1]{(#1 points) \ }
\newcommand{\Point}[1]{(#1 point) \ }


%% For Answers:
\newcounter{answernumber}
\newcommand{\Answer}{%
  \refstepcounter{answernumber}%
  \setcounter{partnumber}{0}%
  \setcounter{subpartnumber}{0}%
  \item[\makebox{\hfil\textbf{\theanswernumber.}\hfil}]%
  }


% Setting up the headers for the first page
%\chead{} % will default to what is typed in the file.
\fancyhead[L]{\textsc{Math 22a}}
\fancyhead[R]{\textsc{Fall 2024}}
\fancyfoot[R]{
	\vspace*{\fill}\footnotesize\textcopyright{} \emph{Harvard Math 22a}	
}
\renewcommand{\headrulewidth}{0pt}

%\fancyhead[C]{}
%	\chead{}	
%	\lhead{}
%	\rhead{}
%	\cfoot{}


% Putting copyright on the footer of each page
%\fancypagestyle{secondpage}{
%	\fancyhead[C]{TESTing again} % change this to be blank
%		%\chead{}	
%	\fancyhead[L]{bears} % change this to be blank
%	\fancyhead[R]{dogs} % change this to be blank
%	\fancyfoot[R]{ %keeps the footer the same
%		\vspace*{\fill}\footnotesize\textcopyright{} \emph{Harvard Math 22a}	
%	}
%}
%\renewcommand{\headrulewidth}{0pt}
%\pagestyle{fancy}
\pagestyle{empty}


%\lhead{}
%\rhead{}
%\rfoot{\vspace*{\fill}\footnotesize\textcopyright{} \emph{Harvard Math 22a}}
%\renewcommand{\headrulewidth}{0pt}
%\pagestyle{fancy}




%%
%% Common shortcut commands:
%%
\newcommand{\ds}{\displaystyle}
\newcommand{\sgn}{\operatorname{sgn}}
\renewcommand{\Pr}{\operatorname{P}}
\newcommand{\CPr}[3][{}]{\Pr_{\text{#1}}\left(#2\,|\,#3\right)}

\newcommand{\C}{\mathbb{C}}
\newcommand{\R}{\mathbb{R}}
\newcommand{\Z}{\mathbb{Z}}
\newcommand{\N}{\mathbb{N}}
\newcommand{\Q}{\mathbb{Q}}
\newcommand{\D}{\mathbb{D}}

\newcommand{\rref}{\operatorname{rref}}
\newcommand{\im}{\operatorname{im}}
\newcommand{\rank}{\operatorname{rank}}
\newcommand{\diag}{\operatorname{diag}}
\newcommand{\Span}{\operatorname{span}}
%\renewcommand{\vec}[1]{\mathbf{#1}}
\newcommand{\charpoly}{\mathscr{P}}
\newcommand{\nicefrac}[2]{\sfrac{#1}{#2}} % using xfrac package
\newcommand{\topology}{\mathcal{T}}
\newcommand{\Inn}{\operatorname{Inn}}

\newcommand{\blankbox}[2]{\fbox{\rule{#1}{0in}\rule{0in}{#2}}}

\newenvironment{myindentpar}[1]%
 {\begin{list}{}%
         {\setlength{\leftmargin}{#1}
          \setlength{\rightmargin}{\leftmargin}}%
         \item[]%
 }
 {\end{list}}
\newtheorem{theorem}{Theorem}
\newtheorem*{NStheorem}{Theorem}
\newtheorem{cor}[theorem]{Corollary}
\newtheorem*{claim}{Claim}
\newtheorem{defn}{Definition}

\newenvironment{amatrix}[1]{%
  \left(\begin{array}{@{}*{#1}{c}|c@{}}
}{%
  \end{array}\right)
}

%--------grstep
% For denoting a Gauss' reduction step.
% Use as: \grstep{\rho_1+\rho_3} or \grstep[2\rho_5 \\ 3\rho_6]{\rho_1+\rho_3}
\newcommand{\grstep}[2][\relax]{%
   \ensuremath{\mathrel{
       {\mathop{\longrightarrow}\limits^{#2\mathstrut}_{
                                     \begin{subarray}{l} #1 \end{subarray}}}}}}
\newcommand{\swap}{\leftrightarrow}


\usetikzlibrary{arrows}
\usepackage{wrapfig}

\makeatletter
\renewcommand*\env@matrix[1][*\c@MaxMatrixCols c]{%
	\hskip -\arraycolsep
	\let\@ifnextchar\new@ifnextchar
	\array{#1}}
\makeatother

\def\env@matrix{\hskip -\arraycolsep
	\let\@ifnextchar\new@ifnextchar
	\array{*\c@MaxMatrixCols c}}

\chead{\Large\textbf{Diagonalization, $\vec\S$5.2}}% 

\begin{document}
\thispagestyle{fancy}

Recall from last class:

\begin{itemize}
\item If $A$ is an $n \times n$ matrix, then $\lambda \in \mathbb{R}$ is an {\bf eigenvalue} of $A$ if there exists a non-zero vector $v$, called an {\bf eigenvector} corresponding to $\lambda$, such that $$Av=\lambda v.$$
\item If $\lambda$ is an eigenvalue of an $n \times n$ matrix $A$, then the set of all solutions $v$ to $Av = \lambda v$ is the {\bf eigenspace} of $A$ corresponding to $\lambda$.
\item {\bf Theorem 1:} The eigenvalues of a triangular matrix are the diagonal entries.
\item {\bf Theorem 2:} If $v_1, \dots, v_p$ are eigenvectors corresponding to distinct eigenvalues $\lambda_1, \dots, \lambda_p$ of an $n\times n$ matrix $A$, then $v_1, \dots, v_p$ are linearly independent.
\item If $A$ is an $n\times n$ matrix, then the characteristic polynomial of $A$, denoted by $p_A(\lambda)$, is given by $$p_A(\lambda)=\text{det}(A-\lambda I_n).$$
\item $\lambda$ is an eigenvalue of an $n\times n$ matrix $A$ if and only if $\lambda$ is a zero of the characteristic polynomial.
\item If $A$ and $B$ are similar matrices, then they have the same characteristic polynomials and hence the same eigenvalues (with the same algebraic multiplicities). 
\end{itemize}

\newpage

\begin{enumerate}

\Problem Let $A=\begin{bmatrix} 1 & 1\\ -2 & 4\\ \end{bmatrix}$. Find eigenvalues and eigenspaces of $A$.

\vfill

\begin{defn}
An $n\times n$ matrix $A$ is {\bf diagonalizable} if $A$ is similar to a diagonal matrix.
\end{defn}

\newpage

{\bf Diagonalization Theorem:} An $n \times n$ matrix $A$ is diagonalizable if and only if $A$ has $n$ linearly independent eigenvectors. Moreover, $A=PDP^{-1}$ with $D$ diagonal if and only if the columns of $P$ are $n$ linearly independent eigenvectors of $A$. In this case the diagonal entires of $D$ are the corresponding eigenvalues of $A$.

\begin{defn} A basis of $\mathbb{R}^n$ diagonalizing $A$ is called an {\bf eigenbasis}. 
\end{defn}

\Problem Prove the Diagonalization Theorem.

\vfill

\Problem Let $A=\begin{bmatrix} 
1 & 3 & 3\\
-3 & -5 & -3\\
3 & 3 & 1\\
\end{bmatrix}$ If possible, diagonalize $A$.

\vfill

\newpage

\Problem Let $A=\begin{bmatrix} 
2 & 4 & 3\\
-4 & -6 & -3\\
3 & 3 & 1\\
\end{bmatrix}$ If possible, diagonalize $A$.


\vfill
\vfill

{\bf Theorem 6:} An $n \times n$ matrix $A$ with $n$ distinct eigenvalues is diagonalizable.

\Problem Prove the theorem.

\vfill

\newpage


\end{enumerate}

\begin{comment}
\newpage
\chead{\Large\textbf{Diagonalization-- Answers / Solutions}}%
\lhead{}
\rhead{}
\thispagestyle{fancy}

\begin{enumerate}
  \Answer In order to find the eigenvalues, we compute the zeros of the characteristic polynomial. Note $p_A(\lambda)= \det \begin{bmatrix} 1-\lambda & 1\\ -2 & 4- \lambda\\ \end{bmatrix}=(1-\lambda)(4-\lambda)+2=(\lambda-2)(\lambda-3)$. Thus the eigenvalues of $A$ are 2 and 3. We find the eigenspaces for each now.
  
\underline{$\lambda=2$}. In this case, we get $$A-2I_2= \begin{bmatrix} -1 & 1\\ -2 & 2\\ \end{bmatrix}\sim \begin{bmatrix} -1 & 1\\ 0 & 0\\ \end{bmatrix}.$$ Thus the eigenspace corresponding to the eigenvalue of 2 is given by $$E_2 = N(A-2I_2) = \text{Span}\left \{\begin{bmatrix} 1\\ 1\\ \end{bmatrix} \right\}.$$

\underline{$\lambda=3$}. In this case, we get $$A-3I_2= \begin{bmatrix} -2 & 1\\ -2 & 1\\ \end{bmatrix}\sim \begin{bmatrix} -2 & 1\\ 0 & 0\\ \end{bmatrix}.$$ Thus the eigenspace corresponding to the eigenvalue of 3 is given by $$E_3 = N(A-3I_2) = \text{Span}\left \{\begin{bmatrix} 1\\ 2\\ \end{bmatrix} \right\}.$$

Thus we get $$A= \begin{bmatrix} 1 & 1\\ 1 & 2\\ \end{bmatrix} \begin{bmatrix} 2 & 0\\ 0 & 3\\ \end{bmatrix} \begin{bmatrix} 1 & 1\\ 1 & 2\\ \end{bmatrix}^{-1}.$$
    \Answer
    \begin{proof} We prove the reverse implication first. 
    
    $\impliedby$. Let $v_1, \dots, v_n$ be linearly independent eigenvectors with corresponding eigenvalues $\lambda_1, \dots, \lambda_n$. Define $P=[v_1 \; \dots \; v_n]$ and $D=\begin{bmatrix} \lambda_1 & 0 & \dots & 0 & 0\\
    0 & \lambda_2 & \dots & 0 & 0\\
    \vdots & \vdots & \ddots & \vdots & \vdots\\
    0 & 0& \dots & 0 & \lambda_n\\
    \end{bmatrix}.$ Then 
    \begin{align*}
    AP &= A[v_1 \; \dots \; v_n ]\\
    & = [Av_1 \; \dots \; Av_n]\\
    &= [\lambda_1 v_1 \; \dots \; \lambda_n v_n]\\
    &= PD.
    \end{align*}
    Since $P$ is invertible, we get $A=PDP^{-1}$, and hence $A$ is diagonalizable. 
    
    $\implies$. Suppose $A=PDP^{-1}$ for some matrix $P$ and some diagonal matrix $D$. Let $P=[v_1 \; \dots \; v_n]$ and let $D$ be diagonal matrix with diagonal entries $\lambda_1, \dots, \lambda_n$. Since $P$ is invertible, the Invertible Matrix Theorem gives $v_1, \dots, v_n$ are linearly independent. Thus we only need to prove that $v_1, \dots, v_n$ are eigenvectors of $A$. Let $e_1, \dots, e_n$ denote the standard basis vectors for $\mathbb{R}^n$, then note $v_k = P e_k$ and hence $P^{-1}v_k = e_k$. Thus
    $$Av_k = PDP^{-1} v_k=PD e_k=P\lambda_k e_k=\lambda_k P e_k =\lambda_k v_k.$$ Thus $v_1,\dots, v_k$ are eigenvectors of $A$ with eigenvalues $\lambda_1, \dots, \lambda_k$.
    \end{proof}
     \Answer In order to find the eigenvalues, we compute the zeros of the characteristic polynomial. Note $p_A(\lambda)= \det \begin{bmatrix} 1-\lambda & 3 & 3\\ -3 & -5 - \lambda & -3\\ 3 & 3 & 1-\lambda \end{bmatrix}=-(\lambda-1)(\lambda-2)^2$. Thus the eigenvalues of $A$ are 1 and -2 (with algebraic multiplicity 2). We find the eigenspaces for each now.
  
\underline{$\lambda=1$}. In this case, we get $$A-I_3= \begin{bmatrix} 0 & 3 & 3\\ -3 & -6 &-3\\ 3 & 3 & 0 \end{bmatrix}\sim \begin{bmatrix} 1 & 0 & -1\\ 0 & 1 & 1\\ 0 & 0 & 0\\ \end{bmatrix}.$$ Thus the eigenspace corresponding to the eigenvalue of 1 is given by $$E_1 = N(A-I_3) = \text{Span}\left \{\begin{bmatrix} 1\\ -1\\ 1\\  \end{bmatrix} \right\}.$$

\underline{$\lambda=-2$}. In this case, we get $$A+2I_3= \begin{bmatrix} 3 & 3 & 3\\ -3 & -3 & -3\\ 3 & 3 & 3\\ \end{bmatrix}\sim \begin{bmatrix} 1 & 1 & 1\\ 0 & 0 & 0\\ 0 & 0 & 0 \\ \end{bmatrix}.$$ Thus the eigenspace corresponding to the eigenvalue -2 is given by $$E_{-2} = N(A+2I_3) = \text{Span}\left \{\begin{bmatrix} -1\\0\\1\\ \end{bmatrix},\begin{bmatrix} -1\\1\\0\\ \end{bmatrix} \right\}.$$

Thus we get $$A= \begin{bmatrix} 1 & -1 &-1\\ -1 & 1 & 0\\ 1 & 0 & 1 \\ \end{bmatrix} \begin{bmatrix} 1 & 0 & 0\\ 0 & -2 & 0\\ 0 & 0 & -2\\ \end{bmatrix} \begin{bmatrix} 1 & -1 &-1\\ -1 & 1 & 0\\ 1 & 0 & 1 \\ \end{bmatrix}^{-1}.$$

\Answer In order to find the eigenvalues, we compute the zeros of the characteristic polynomial. Note $p_A(\lambda)= \det \begin{bmatrix} 2-\lambda & 4 & 3\\ -4 & -6 - \lambda & -3\\ 3 & 3 & 1-\lambda \end{bmatrix}=-(\lambda-1)(\lambda-2)^2$. Thus the eigenvalues of $A$ are 1 and -2 (with algebraic multiplicity 2). We find the eigenspaces for each now.
  
\underline{$\lambda=1$}. In this case, we get $$E_1 = N(A-I_3) = \text{Span}\left \{\begin{bmatrix} 1\\ -1\\ 1\\  \end{bmatrix} \right\}.$$

\underline{$\lambda=-2$}. In this case, we get $$E_{-2} = N(A+2I_3) = \text{Span}\left \{\begin{bmatrix} -1\\1\\0\\ \end{bmatrix} \right\}.$$

Hence we have only two linearly independent eigenvectors and by the Diagonalization theorem, $A$ cannot be diagonalized.


%\Answer 
%
%\begin{proof}
%Suppose $A$ is an $n\times n$ matrix with $n$ distinct eigenvalues. Let $v_1, \dots, v_n$ be $n$ eigenvectors each corresponding to a different eigenvalue. We proved last time that eigenvectors corresponding to different eigenvalues must be linearly independent. Thus by the Diagonalization Theorem, $A$ must be diagonalizable. 
%\end{proof}
%\Answer 			A vector $p(t) = a_0 + a_1 t + a_2 t^2 + a_3 t^3$ will be in the kernel if $a_0 - a_1 = a_2 - a_3 = 0$.
%				Thus a basis for the kernel will be 
%					$\left\{ 1+t, t^2+t^3 \right\}$.
%					
%					Let $\mathcal{B} = \{ 1,t,t^2,t^3\}$ be a basis for $\mathbb{P}_3$, and $\mathcal{C} = \left\{ \begin{bmatrix} 1 & 0 \\ 0 & 0 \end{bmatrix},\begin{bmatrix} 0 & 1 \\ 0 & 0 \end{bmatrix}, \begin{bmatrix} 0 & 0 \\ 1 & 0 \end{bmatrix}, \begin{bmatrix} 0 & 0 \\ 0 & 1 \end{bmatrix} \right\}$
%			be a basis for $M_{2 \times 2}$.
%			
%							We need to find the matrix
%					$\left[ \begin{array}{ccc} [T(b_1)]_{\mathcal{C}} & \cdots & [T(b_n)]_{\mathcal{C}} \end{array} \right]$:				\begin{align*}
%					T(b_1) &= T(1) = \begin{bmatrix} 1 & 0 \\ -1 & 0 \end{bmatrix} \\
%					T(b_2) &= T(t) = \begin{bmatrix} -1 & 0 \\ 1 & 0 \end{bmatrix} \\
%					T(b_3) &= T(t^2) = \begin{bmatrix} 0 & 1 \\ 0 & -1 \end{bmatrix} \\
%					T(b_4) &= T(t^3) = \begin{bmatrix} 0 & -1 \\ 0 & 1 \end{bmatrix}
%				\end{align*}
%				Thus the matrix is
%				\[
%					\begin{bmatrix} 1 & -1 & 0 & 0 \\ 0 & 0 & 1 & -1 \\ -1 & 1 & 0 & 0 \\ 0 & 0 & -1 & 1 \end{bmatrix}.
%				\]
%				
%		From the previous part, we know that the range of $T$ is the span of 
%				\[
%					\left\{ \begin{bmatrix} 1 & 0 \\ -1 & 0 \end{bmatrix}, \begin{bmatrix} -1 & 0 \\ 1 & 0 \end{bmatrix}, \begin{bmatrix} 0 & 1 \\ 0 & -1 \end{bmatrix}, \begin{bmatrix} 0 & -1 \\ 0 & 1 \end{bmatrix} \right\}.
%				\]
%
%				Removing the dependent entries, we find a basis
%				\[
%					\mathcal{D} = \left\{ \begin{bmatrix} 1 & 0 \\ -1 & 0 \end{bmatrix}, \begin{bmatrix} 0 & 1 \\ 0 & -1 \end{bmatrix} \right\}
%				\]
%					
%					
%\Answer Evaluating $T$ on the standard basis for $\mathbb{P}_2$ yields: 
%$$T\left(1 \right) = \begin{bmatrix} 1 & 0\\ 0 & 0\end{bmatrix};  \; \;
%T\left( t \right) = \begin{bmatrix} 0 & 1\\ 0 & 0\end{bmatrix}; \; \;
%T\left( t^2 \right) = \begin{bmatrix} 0 & 2\\ 0 & 2\end{bmatrix}.$$
%
%Let $\mathcal{E}$ denote the standard basis of $U_{2 \times 2}$, then 
%$$[T\left(1 \right)]_{\mathcal{E}} = \begin{bmatrix} 1\\ 0\\  0\\\end{bmatrix};  \; \;
%[T\left( t \right)]_{\mathcal{E}} = \begin{bmatrix} 0 \\ 1\\ 0\\\end{bmatrix}; \; \;
%[T\left( t^2 \right)]_{\mathcal{E}} = \begin{bmatrix} 0 \\ 2 \\ 2\end{bmatrix}.$$
%
%Thus the matrix of $T$ relative to the standard coordinates is given by
%$$M=\begin{bmatrix}
%1 & 0 & 0\\
%0 & 1 & 2\\
%0 & 0 & 2\\
%\end{bmatrix}.$$
%
%We proved in class that the eigenvalues of an upper triangular matrix are given by the diagonal entries, and hence the eigenvalues of $T$ are $1$ and $2$.
%
%We first find the eigenspace corresponding to the 1-eigenvalue.  Thus, we have
%$$E_1 = N(M- I)= \text{Span}\left\{  \begin{bmatrix} 1\\0\\0\\\end{bmatrix}, \begin{bmatrix} 0\\1\\0\\\end{bmatrix}\right\}.$$
%
%Similarly, we find the eigenspace for $\lambda=2$.  That is, we have
%$$E_2 = N(M-2I)= \text{Span}\left\{  \begin{bmatrix} 0\\2\\1\\ \end{bmatrix}\right\}.$$
%
%Since we have 3 linearly independent eigenvectors, the diagonalization theorem implies $T$ will act diagonally in the bases coming from the eigenvectors.  Namely, we have the eigenbasis $$\left\{  \begin{bmatrix} 1\\0\\0\\\end{bmatrix}, \begin{bmatrix} 0\\1\\0\\\end{bmatrix}, \begin{bmatrix} 0\\2\\1\\ \end{bmatrix}\right\}.$$  
%
%This gives the corresponding bases $\mathcal{P} =\left\{1, t, 2t+t^2 \right \}$ in $\mathbb{P}_2$ and $ \mathcal{B}= \left\{  \begin{bmatrix} 1 & 0\\ 0 & 0\end{bmatrix}, \begin{bmatrix} 0 & 1\\ 0 & 0\end{bmatrix}, \begin{bmatrix} 0 & 2\\ 0 & 1\end{bmatrix}\right\}$ for $U_{2 \times 2}$.  Thus the matrix $M$ of $T$ relative to $\mathcal{P}$ and $\mathcal{B}$ is given by
%$$M=\begin{bmatrix}
%1 & 0 & 0\\
%0 & 1 & 0\\
%0 & 0 & 2\\
%\end{bmatrix}.$$
%\Answer 
%\begin{enumerate}
%\item The kernel of $T$ is the nullspace of $A$, so we row reduce $A$ to
%find 
%$$\begin{bmatrix} 1 & 0 & -1 \\ 0 & 1 & 0 \\ 0 & 0 & 0\end{bmatrix}.$$
%Hence a basis for the kernel consists of the one vector
%$$\begin{bmatrix}1\\0\\1\end{bmatrix},$$ 
%which means that the kernel of $T$ is one-dimensional.
%\item We already row reduced $A$ and we found two pivots, which means that
%the rank of $A$ is $2$.
%
%The image of $T$ is the column space of $A$ and a basis for that space
%consists of the first two columns of $A$:
%$$\begin{bmatrix}0\\0\\-1\end{bmatrix}, \begin{bmatrix}1\\1\\1\end{bmatrix}.$$
%\end{enumerate}
%\Answer
%\begin{enumerate}
%\item The characteristic equation of $A$ is
%$$\det(A - \lambda I) = \begin{vmatrix}-10 - \lambda & -18 \\ 9 & 17 -
%  \lambda\end{vmatrix}
%= \lambda^2 - 7\lambda - 8.$$
%The roots of this equation are $-1$ and $8$, which are therefore the
%eigenvalues of $A$.
%\item For the eigenvalue $-1$, the eigenspace is the nullspace
%of $$\begin{bmatrix}-9&-18\\9&18\end{bmatrix} \sim \begin{bmatrix} 1 &
%  2 \\ 0 & 0 \end{bmatrix}.$$
%Therefore a basis for this eigenspace consists of the vector
%$\begin{bmatrix}-2\\1\end{bmatrix}$, for example.
%
%For the eigenvalue $8$, the eigenspace is the nullspace
%of $$\begin{bmatrix}-18&-18\\9&9\end{bmatrix} \sim \begin{bmatrix} 1 &
%  1 \\ 0 & 0 \end{bmatrix}.$$
%Therefore a basis for this eigenspace consists of the vector
%$\begin{bmatrix}-1\\1\end{bmatrix}$, for example.
%\item We know that $A = PDP^{-1}$ where $D
%= \begin{bmatrix}-1&0\\0&8\end{bmatrix}$ and $P = \begin{bmatrix} -2 &
%  -1 \\ 1 & 1\end{bmatrix}$.
%Consider the diagonal matrix $E = \begin{bmatrix} -1 & 0 \\ 0 &
%  2 \end{bmatrix}$. 
%Note that $E^3 = D$ and hence $B^3 = A$ where $B = PEP^{-1}$.
%Thus, the matrix 
%$$B = \begin{bmatrix}-2 & -1 \\ 1 & 1 \end{bmatrix}\begin{bmatrix}-1 &
%  0 \\ 0 & 2 \end{bmatrix} \begin{bmatrix} -1 & -1 \\ 1 &
%  2 \end{bmatrix} = \begin{bmatrix} -4 & -6 \\ 3 & 5 \end{bmatrix}$$ 
%works!
%\end{enumerate}
%\Answer
%\begin{enumerate}
%\item We must show that $T((p+q)(t)) = T(p(t)) + T(q(t))$ and $T((cp)(t)) =
%cT(p(t))$ hold for all polynomials $p(t)$ and $q(t)$ in $\mathbb{P}_2$
%and every scalar $c \in \mathbb{R}$.
%
%Both equalities are immediate since $(p+q)(t)$ is defined to equal
%$p(t) + q(t)$ and $(cp)(t)$ is defined to equal $c(p(t))$ for every
%value of $t$.
%\item Write $p(t) = a_2t^2 + a_1t + a_0$ so that $[p(t)]_{\mathcal{B}}
%= \begin{bmatrix}a_0\\a_1\\a_2\end{bmatrix}$. 
%Then $$\begin{bmatrix}p(1)\\p(2)\\p(3)\end{bmatrix}
%= \begin{bmatrix}a_2 + a_1 + a_0\\a_24 + a_12 + a_0\\a_29 + a_13 +
%  a_0\end{bmatrix}
%= \begin{bmatrix}1&1&1\\1&2&4\\1&3&9\end{bmatrix}\begin{bmatrix}a_0\\a_1\\a_2\end{bmatrix},$$
%which shows that $M = \begin{bmatrix}1&1&1\\1&2&4\\1&3&9\end{bmatrix}$ works.
%\item We are asked to solve the equation $T(p(t))
%= \begin{bmatrix}c_1\\c_2\\c_3\end{bmatrix}$ for the polynomial $p(t)
%= a_2t^2 + a_1t + a_0$.
%By part~(b), this equation is equivalent to the matrix
%equation $$\begin{bmatrix}c_1\\c_2\\c_3\end{bmatrix}
%= \begin{bmatrix}1&1&1\\1&2&4\\1&3&9\end{bmatrix}\begin{bmatrix}a_0\\a_1\\a_2\end{bmatrix}.$$
%This equation always has a solution since the matrix $M$ is invertible.
%\end{enumerate}
%
%\Answer Let $P = \begin{bmatrix} \mathbf{b}_1 & \mathbf{b}_2 &
%  \mathbf{b}_3 \end{bmatrix}$. Then $$A = PDP^{-1}, \qquad B =
%PEP^{-1}$$ where $$D = \begin{bmatrix} \lambda_1 & 0 & 0 \\ 0 &
%  \lambda_2 & 0 \\ 0 & 0 & \lambda_3\end{bmatrix}, \qquad E
%= \begin{bmatrix} \mu_1 & 0 & 0 \\ 0 & \mu_2 & 0 \\ 0 & 0 &
%  \mu_3 \end{bmatrix}.$$
%
%Then $$AB = (PDP^{-1})(PEP^{-1}) = PDEP^{-1} = PEDP^{-1}
%=(PDP^{-1})(PEP^{-1}) = BA,$$ since $$DE = \begin{bmatrix} \lambda_1\mu_1 & 0 & 0 \\ 0 &
%  \lambda_2\mu_2 & 0 \\ 0 & 0 & \lambda_3\mu_3\end{bmatrix} = ED.$$
%\Answer
%\begin{enumerate}
%\item We need to show that $K_j \subseteq K_{j+1}$.
%So suppose $x \in K_j$, that is $A^jx = 0$. Then $$A^{j+1}x
%=A(A^jx) = A\mathbf0 = \mathbf0,$$ which shows that $x \in K_{j+1}$.
%\item Suppose that $K_r = K_{r+1}$. We prove by induction on $n \geq 1$ that
%$K_r = K_{r+n}.$
%
%The base case $n = 1$ is given to us!
%
%For the induction step, suppose we have shown that $K_r = K_{r+n}$.
%To show that $K_r = K_{r+n+1}$ we can show instead that $K_{r+n} =
%K_{r+n+1}$.
%We already know that $K_{r+n} \subseteq K_{r+n+1}$ so it suffices to
%show that $K_{r+n+1} \subseteq K_{r+n}$.
%
%Suppose that $x \in K_{r+n+1}$, that is $A^{r+n+1}x = \mathbf0$.
%Then $A^{r+1}(A^n x) = \mathbf0$, which means that $A^n x \in
%K_{r+1}$. Since $K_{r+1} = K_r$, we conclude that $A^n x \in K_r$ and
%hence $A^r(A^n x) = \mathbf0$. Since $A^rA^n = A^{n+r}$, we conclude
%that $x \in K_{n+r}$.
%\end{enumerate}
%\begin{enumerate}
%\item True. An example of such a matrix is $$A
%= \begin{bmatrix}1&-1&-1&-1\\0&0&0&0\end{bmatrix}.$$
%Indeed, a basis for the eigenspace of this matrix consists of the
%three
%vectors $$\begin{bmatrix}1\\0\\0\\1\end{bmatrix},\begin{bmatrix}1\\0\\1\\0\end{bmatrix},\begin{bmatrix}1\\1\\0\\0\end{bmatrix},$$
%which is also a basis of 
%$$\mathrm{Span}\left\{\begin{bmatrix}1\\0\\0\\1\end{bmatrix},\begin{bmatrix}1\\0\\1\\0\end{bmatrix},\begin{bmatrix}1\\1\\0\\0\end{bmatrix},\begin{bmatrix}3\\1\\1\\1\end{bmatrix}\right\}$$
%since the last spanning vector is redundant.
%\item True.  We need to show that $\{ v_1, \cdots, v_k \}$ is linearly independent and spans $V$.  Let $v \in V$, then $T(v) \in W$.  Since $\{ T(v_1), \cdots, T(v_k)\}$ is a basis for $W$, there exists $c_1, \cdots, c_k \in \mathbb{R}$ so that $$T(v) = c_1 T(v_1) +\cdots T(v_k).$$  Since $T$ is injective and linear, we get $v = c_1 v_1+\cdots +c_k v_k$.  Thus, $v_1, \cdots, v_k$ span $V$.
%
%We are left to check linear independence.  Consider the vector equation $d_1 v_1 + \cdots d_k v_k = 0_V$.  Applying $T$ to both sides and using the linearity of $T$ yields $$d_1 T(v_1) + \cdots d_k T(v_k) = 0.$$  Since $T(v_1), \cdots, T(v_k)$ are linearly independent, we know that $d_1=d_2 = \cdots = d_k =0$.  Thus $v_1, \cdots, v_k$ are linearly independent, and hence they form a basis. 
%\item False.  The zero matrix is not invertible, so the set cannot be a subspace.
%\item True.  Since the sum of each row of $A$ equals 1, we know that 1 is an eigenvalue of $A$ with eigenvector of all 1's.  Thus $A-I$ must have non-trivial null space, and hence by the invertible matrix theorem, $A-I_n$ is not invertible.
%\item True.  Let $v \in E_1 \cap E_2$.  Let $E_j$ be the eigenspace corresponding to $\lambda_j$ for $j=1,2$.  Then $\lambda_1 v =\lambda_2 v$ since $\lambda_1\neq \lambda_2$, we must have $v=0$.
%\item True. Let $x_1,x_2,x_3$ be a basis for $W_1$ and $y_1, y_2, y_3$ be a basis for $W_2$. Since $x_1, x_2, x_3, y_1, y_2, y_3$ is 6 vectors in $\mathbb{R}^5$, the vectors must be linearly dependent. Thus there are scalars, not all zero, such that $$c_1 x_1 +c_2 x_2+ c_3 +c_4 y_1+ c_5 y_2+c_6 y_3 =0,$$ and hence  $$c_1 x_1 +c_2 x_2 + c_3 x_3 = -c_4 y_1-c_5 y_2-c_6 y_3.$$ Thus $c_1 x_1 +c_2 x_2 + c_3 x_3$ is a non-zero element in $W_1 \cap W_2$.
%\item False. Let $V=\text{Span}\left\{\begin{bmatrix} 1\\ 0\\ \end{bmatrix} \right\}$ and $W=\text{Span}\left\{\begin{bmatrix} 0\\ 1\\ \end{bmatrix} \right\}$. Now $\begin{bmatrix} 1\\ 0\\ \end{bmatrix}, \begin{bmatrix} 0\\ 1\\ \end{bmatrix} \in V\cup W$, but $\begin{bmatrix} 1\\ 0\\ \end{bmatrix}+ \begin{bmatrix} 0\\ 1\\ \end{bmatrix}=\begin{bmatrix} 1\\ 1\\ \end{bmatrix}\not \in V \cup W.$
%\item False. Let $W_1=\text{Span}\left\{\begin{bmatrix} 1\\ 0\\ \end{bmatrix} \right\}$ and $W_2=\text{Span}\left\{\begin{bmatrix} 0\\ 1\\ \end{bmatrix} \right\}.$ Then $\dim W_1 \leq \dim W_2$, but $W_1$ is not a subset of $W_2$.
%\item True. Suppose that $Ax = \lambda x$. Then $$A^3x = A^2(Ax) =
%A^2(\lambda x) = \lambda A^2 x = \lambda A(A x) = \lambda A(\lambda x)
%= \lambda^2 A x = \lambda^3 x,$$ which shows that $x$ is an
%eigenvector of $A^3$ with eigenvalue $\lambda^3$.
%\item False. Take $A = \begin{bmatrix}1&0\\0&-1\end{bmatrix}$. Then $x
%= \begin{bmatrix}1\\0\end{bmatrix}$ and $y
%= \begin{bmatrix}0\\1\end{bmatrix}$ are both eigenvectors of $A$, but
%$x+y$ is not an eigenvector of
%$A$ since
%$$A(x+y)
%= \begin{bmatrix}1&0\\0&-1\end{bmatrix}\begin{bmatrix}1\\1\end{bmatrix}
%= \begin{bmatrix}1\\-1\end{bmatrix}$$ is not a multiple of $x + y$.
%\item True. Because $T$ is one-to-one, the dimension of the image of $T$
%must be $n$. Since $W$ has dimension $n$, the image of $T$ must equal $W$.
%%\item True. The characteristic polynomial has three distinct roots: $0$,
%%$2$, $5$. Therefore $A$ and $B$ are both similar to the diagonal
%%matrix $\begin{bmatrix}0&0&0\\0&2&0\\0&0&5\end{bmatrix}$, which means
%%that they are also similar to eachother.
%\end{enumerate}
        
          
\end{enumerate}  
\end{comment}

\end{document}


%%% Local ables: 
%%% mode: latex
%%% TeX-master: t
%%% End: 
